% Vietnamese translation for OI Public Library
% Source-File: IOI中国国家候选队论文/2009/刘聪/浅谈数位类统计问题.ppt
\documentclass[11pt]{article}
\usepackage{oipl-vn}

\newcommand{\Oh}{\mathcal{O}}

\title{Bàn sơ lược về các bài toán thống kê theo chữ số\\{\large Ghi chú theo slide}}
\author{Lưu Thông\\Trường Trung học số 2 Thanh Đảo, Sơn Đông}
\date{2009}

\begin{document}
\maketitle

\origtitle{On digit-statistics problems}
\sourcepath{IOI-China-Candidate-Team-Papers/2009/Liu-Cong/digit-statistics-slides.ppt}

\section{Mạch trình bày}

Slide dùng năm ví dụ để dẫn người đọc từ đếm chữ số cơ bản đến các bài có
thứ tự từ điển và nhiều ràng buộc hơn. Điểm chung là không duyệt từng số, mà
phân tích biểu diễn theo cơ số rồi đếm theo tiền tố.

\begin{enumerate}
  \item \textit{Amount of degrees}: đếm số có đúng \(k\) chữ số \(1\) trong
  biểu diễn cơ số \(B\).
  \item \textit{Sorted bit sequence}: xác định phần tử theo thứ tự sau khi
  nhóm bởi số bit \(1\).
  \item \textit{Sequence}: gắn truy hồi đếm với vị trí trong dãy sinh ra bởi
  quy luật chữ số.
  \item \textit{Tickets}: dùng quy hoạch động trên tổng chữ số hai nửa vé.
  \item \textit{Graduated Lexicographical Ordering}: kết hợp đếm theo độ dài,
  trọng số chữ cái và thứ tự từ điển.
\end{enumerate}

\section{Ý tưởng kỹ thuật}

Với một giới hạn \(N\), cách làm quen thuộc là quét các chữ số của \(N\) từ
cao xuống thấp. Khi đặt một chữ số nhỏ hơn chữ số hiện tại của \(N\), phần
đuôi còn lại được đếm bằng tổ hợp hoặc quy hoạch động; khi đặt bằng, ta tiếp
tục xét chữ số kế tiếp. Nhờ vậy truy vấn dạng ``có bao nhiêu số không vượt quá
\(N\)'' thường giảm về \(\Oh(\log N)\) trạng thái lớn.

Trong các ví dụ có yêu cầu tìm phần tử thứ \(K\), thao tác đếm được dùng như
một thủ tục kiểm tra. Ta thử từng lựa chọn cho chữ số kế tiếp, trừ đi số lượng
phương án bị bỏ qua, rồi cố định lựa chọn đầu tiên làm cho thứ hạng rơi vào
khối đang xét.

\section{Ghi chú khi đọc cùng bài viết}

Bài viết đầy đủ trình bày đề gốc và các chi tiết chứng minh. Bản slide nên
được đọc như bản đồ thuật toán: nhận ra đại lượng cần đếm, chọn trạng thái cho
phần tiền tố đã cố định, rồi dùng công thức đếm phần hậu tố để tránh duyệt
tuyến tính trên miền số.

\end{document}
